\section*{Bab \ref{ch:g9:thales} --- Teorema Thales}

\begin{solution}{exo:g9:thales:1}
Thales: $\dfrac{AM}{AB} = \dfrac{AN}{AC} = \dfrac{MN}{BC}$, yaitu
$\dfrac26 = \dfrac13$. Jadi $\dfrac{3}{AC} = \dfrac13$ memberi
$AC = 9$, dan $\dfrac{MN}{9} = \dfrac13$ memberi $MN = 3$.
\end{solution}

\begin{solution}{exo:g9:thales:2}
Carilah lebih dahulu $AB = AM + MB = 5 + 3 = 8$. Lalu
$\dfrac{AM}{AB} = \dfrac{AN}{AC}$ memberi $\dfrac58 = \dfrac{4}{AC}$,
jadi $5 \times AC = 32$ dan $AC = 6.4$.
\end{solution}

\begin{solution}{exo:g9:thales:3}
Susunan kupu-kupunya bekerja persis seperti susunan bersarangnya:
$\dfrac{AM}{AB} = \dfrac{AN}{AC} = \dfrac{MN}{BC}$, yaitu
$\dfrac{3}{7.5} = \dfrac25$. Jadi $\dfrac{2}{AC} = \dfrac25$ memberi
$AC = 5$, dan $\dfrac{2.6}{BC} = \dfrac25$ memberi
$BC = \dfrac{2.6 \times 5}{2} = 6.5$.
\end{solution}

\begin{solution}{exo:g9:thales:4}
$\dfrac{AM}{AB} = \dfrac68 = \dfrac34$ dan
$\dfrac{AN}{AC} = \dfrac{9}{12} = \dfrac34$: perbandingannya sama,
titiknya berurutan sama, jadi $(MN) \parallel (BC)$ menurut konvers Thales.
\end{solution}

\begin{solution}{exo:g9:thales:5}
$\dfrac{AM}{AB} = \dfrac46 = \dfrac23$ dan
$\dfrac{AN}{AC} = \dfrac58$. Pemeriksaan silangnya:
$\dfrac23 = \dfrac{16}{24}$ dan $\dfrac58 = \dfrac{15}{24}$:
perbandingannya berbeda. Kalau garisnya \omterm{def:g6:lines:perp}{sejajar}, teorema Thales akan
memaksa keduanya sama --- jadi $(MN)$ dan $(BC)$ \emph{tidak} \omterm{def:g6:lines:perp}{sejajar}.
\end{solution}

\begin{solution}{exo:g9:thales:6}
Misalkan $S$ ujung bayangannya, $P$ puncak kepala orangnya, $L$
puncak lampunya. Orangnya ($1.8$ m pada jarak $3$ m dari $S$) dan
lampunya (tinggi $h$ pada jarak $3 + 2 = 5$ m dari $S$) adalah dua
\omterm{def:g6:lines:objects}{ruas garis} tegak yang \omterm{def:g6:lines:perp}{sejajar} yang dipotong sinar cahaya $(SL)$:
Thales dari titik $S$ memberi
\[
\frac{1.8}{h} = \frac{3}{5},
\qquad\text{jadi}\quad
h = \frac{1.8 \times 5}{3} = 3 \text{ m}.
\]
\end{solution}

\begin{solution}{exo:g9:thales:7}
\emph{1.} Jarak yang sebenarnya: $6.8 \times 25\,000 = 170\,000$ cm
$= 1.7$ km.

\emph{2.} Luas menskala dengan $k^2$: luas yang sebenarnya
$= 8 \times 25\,000^2$
cm$^2$ $= 8 \times 6.25 \times 10^8$ cm$^2$ $= 5 \times 10^9$ cm$^2$.
Karena $1$ km$^2 = 10^{10}$ cm$^2$, itu berarti $0.5$ km$^2$.
\end{solution}

\begin{solution}{exo:g9:thales:8}
\emph{1.} Perbandingannya adalah $\dfrac{AM}{AB} = \dfrac{8}{12} =
\dfrac23$. Thales: $AN = \frac23 \times 15 = 10$ dan
$MN = \frac23 \times 18 = 12$.

\emph{2.} \omterm{def:g9:thales:scaling}{Faktor skalanya} $k = \frac23$. \omterm{def:g6:measure:perimeter}{Kelilingnya}: $ABC$ punya
$12 + 15 + 18 = 45$, dan $AMN$ punya $8 + 10 + 12 = 30 = \frac23 \times
45$: jadi \omterm{def:g6:measure:perimeter}{kelilingnya} pun menskala dengan $k$, seperti setiap panjang.
\end{solution}

\begin{solution}{exo:g9:thales:9}
Di sekitar $O$, garis $(AC)$ dan $(BD)$ berpotongan, dan
$(AB) \parallel (CD)$: yaitu susunan kupu-kupu dengan segitiga $OAB$
dan $OCD$. Thales memberi
\[
\frac{OA}{OC} = \frac{OB}{OD} = \frac{AB}{CD} = \frac46 = \frac23 .
\]
Jadi $\dfrac{OA}{OC} = \dfrac23$, dan kesamaan
$\dfrac{OA}{OC} = \dfrac{OB}{OD}$ mengatakan persis bahwa $O$ memotong
kedua diagonalnya dengan perbandingan yang sama.
\end{solution}

\begin{solution}{pb:g9:thales:1}
\textbf{1.} Pelukisannya menghasilkan dua titik pada $[AB]$;
sebutlah keduanya $M_1$ (dari $P_1$) dan $M_2$ (dari $P_2$).

\textbf{2.} Pada segitiga $A P_3 B$, \omterm{def:g6:lines:perp}{garis sejajar}
$(P_3 B)$ yang melalui $P_1$ dan $P_2$ memotong sisinya
secara \omterm{def:g9:linfunc:linear}{sebanding} (\cref{thm:g9:thales:direct}):
\[
\frac{A M_1}{AB} = \frac{A P_1}{A P_3} = \frac13,
\qquad
\frac{A M_2}{AB} = \frac{A P_2}{A P_3} = \frac23 .
\]
Jangkanya membuat $A P_1$, $P_1 P_2$, $P_2 P_3$ sama, jadi
perbandingannya persis sepertiga --- berapa pun sudut sinarnya dan
bukaan jangka yang dipilih.

\textbf{3.} Untuk seperlima: langkahkan lima panjang yang sama,
hubungkan titik kelimanya ke $B$, lalu gambarlah empat \omterm{def:g6:lines:perp}{garis sejajar}.
Untuk $\frac35$ sebuah \omterm{def:g6:lines:objects}{ruas garis}: gambar yang sama, lalu ambillah
titik yang dipotong \omterm{def:g6:lines:perp}{garis sejajar}
yang melalui $P_3$: titik itu duduk pada $\frac{3}{5}$ dari $[AB]$ dari $A$.

\textbf{4.} $\frac{AM}{MB} = \frac23$ berarti
$AM = \frac25 AB$: lima langkah yang sama pada sinarnya, hubungkan
$P_5$ ke $B$, lalu \omterm{def:g6:lines:perp}{garis sejajar} yang melalui $P_2$ menandai $M$.

\textbf{5.} Mengukur $\sqrt2 = 1.41421\ldots$ cm selalu hanya memakai
desimal yang berhingga banyaknya: setiap ``sepertiga'' yang terukur
adalah hampiran. Pelukisan Thales tidak pernah membaca sebuah
bilangan: ia menyerahkan titik yang persis pada $\frac{\sqrt2}{3}$ ---
yaitu bagian yang sama dari \omterm{def:g6:lines:objects}{ruas garis} yang bahkan tidak dapat
dinyatakan penggaris mana pun (\cref{pb:g9:sqrt:1}).

\textbf{6.} Puncak bendanya, lubangnya dan dasar bayangannya
segaris, dan begitu pula dasar bendanya dan
puncak bayangannya: yaitu dua garis yang berpotongan di lubangnya,
dengan bendanya dan dinding bayangannya \omterm{def:g6:lines:perp}{sejajar}. Thales pada
kupu-kupunya:
\[
\frac{h}{H} = \frac{d}{D},
\qquad\text{jadi}\qquad
h = H \times \frac dD .
\]

\textbf{7.} $h = 12 \times \frac{0.20}{30} = 0.08$ m $= 8$ cm
--- dan terbalik (sinarnya berpotongan di lubangnya).

\textbf{8.} $H = h \times \frac Dd = 9 \times 10^{-3} \times
\frac{1.5 \times 10^{11}}{1} = 1.35 \times 10^{9}$ m. Nilai
sebenarnya $1.39 \times 10^9$ m: sebuah kotak sepatu mengukur
Matahari sampai \omterm{ex:g1:subtraction:difference}{selisih} beberapa persen saja.

\textbf{9.} Bulan: $\frac{3.5 \times 10^6}{3.8 \times 10^8}
\approx 0.0092$. Matahari: $\frac{1.39 \times 10^9}{1.5 \times
10^{11}} \approx 0.0093$. Kedua perbandingannya --- yaitu besar
tampaknya di langit --- hampir sama: jadi Bulan dapat menutupi
Matahari \emph{dengan tepat}, tepi bertemu tepi. Kebetulan kosmis itu
adalah gerhana matahari total.

\textbf{10.} Setiap sinar harus melewati lubang yang satu itu, jadi
sinar dari puncak bendanya berlanjut \emph{ke bawah} dan sinar dari
dasarnya berlanjut \emph{ke atas}: jadi bayangannya terbalik.
Memperbesar lubangnya meloloskan seberkas penuh salinan bayangannya
yang sedikit bergeser, yang saling tumpang tindih: lebih terang,
tetapi mengabur.

\textbf{11.} $\frac{PQ}{MN} = \frac{OP}{OM} = \frac{7.5}{3}
= 2.5$, jadi $PQ = 10$; dan $\frac{ON}{OQ} = \frac{OM}{OP} =
\frac{1}{2.5}$, jadi $ON = \frac{6}{2.5} = 2.4$.

\textbf{12.} Kalau \omterm{def:g6:lines:perp}{garis sejajar} itu melalui titik tengah salah satu
sisinya, perbandingan Thalesnya $\frac12$, jadi garis itu memotong
sisi keduanya di titik tengah \emph{sisi itu}, dan ruas \omterm{def:g6:lines:perp}{garis
sejajarnya} berukuran setengah alasnya: yaitu persis
\cref{thm:g8:midpoints:direct}. Thales adalah teorema titik tengah
yang dibebaskan dari perbandingan $\frac12$.

\textbf{13.} Pada segitiga $ACD$, garis $(EO)$ \omterm{def:g6:lines:perp}{sejajar}
dengan alasnya $(CD)$, dengan $\frac{AO}{AC} = \frac{OA}{OA + OC} =
\frac{4}{4 + 6} = \frac25$. Thales:
$\frac{EO}{DC} = \frac{AO}{AC} = \frac25$, jadi
$EO = \frac25 \times 6 = 2.4$.

\textbf{14.} Pada segitiga $BDC$, dengan cara yang sama
$\frac{BO}{BD} = \frac25$ dan $OF = \frac25 \times 6 = 2.4$.
Karena itu
\[
EF = EO + OF = 4.8
= \frac{2 \times 4 \times 6}{4 + 6} :
\]
yaitu \omterm{pb:g8:speed:1}{rata-rata harmonik} kedua alasnya --- rata-rata perjalanan pulang
pergi (\cref{pb:g8:speed:1}), yang tergambar dalam sebuah trapesium.

\textbf{15.} Aritmetika: $5$; geometri:
$\sqrt{24} \approx 4.90$; harmonik: $4.8$. Urutannya:
$4.8 < 4.90\ldots < 5$, yaitu harmonik $<$ geometri $<$
aritmetika --- dengan kesamaannya hanya untuk alas yang sama. Pada
gambarnya: garis tengahnya berukuran $5$, \omterm{def:g6:lines:perp}{garis sejajar} yang melalui
potongan diagonalnya $4.8$, dan potongan kesebangunannya
$\sqrt{24}$, semuanya bertumpuk di antara alasnya dengan urutan itu.
Ketaksamaannya sudah dibuktikan pada \cref{pb:g8:cosine:1} (geometri
$\leq$ aritmetika, lewat lingkaran dan lewat kesamaan) dan
\cref{pb:g8:speed:1} (harmonik $\leq$ aritmetika); harmonik
$\leq$ geometri menyusul dengan menggabungkan keduanya
($H \times A = G^2$, yang dibuktikan di sana juga).

\end{solution}
